\section{Síntesis y ampliación}

\subsection{Una tabla que recoge toda la sesión}

Esta sesión parte de poker search y culmina en un agent completo para Diplomacy. La tabla siguiente reúne, para cada módulo, el problema que resuelve, su mecanismo principal y sus modos de falla, para que pueda reutilizarse al leer después artículos sobre sistemas multiagente o sobre reasoning.

\begin{longtable}{L{0.20\textwidth}L{0.30\textwidth}L{0.40\textwidth}}
\toprule
Módulo & Problema que resuelve & Mecanismo y riesgos principales \\
\midrule
Base policy & Proponer rápidamente candidatos razonables & Amplia cobertura, pero puede ser tosca en situaciones críticas; no sustituye a search \\
Inference-time search & Mejorar la decisión actual con cómputo adicional & Expansión, evaluación y selección; puede amplificar un evaluator erróneo \\
Human anchor policy & Mantener convention comprensibles para las personas & La imitation aporta un prior; también hereda las debilidades del humano promedio \\
piKL & Equilibrar valor y human-likeness & Una KL penalty limita la deriva de la política; $\lambda$ determina el compromiso \\
Intent model & Convertir el plan en una condición controlable de generación & Las action de corto plazo son claras y verificables; los objetivos de largo plazo se expresan de forma insuficiente \\
Dialogue model & Lograr grounded communication & Genera candidatos condicionados por state, history e intent \\
Action predictor & Predecir el comportamiento de los demás tras recibir un mensaje & Conecta lenguaje y estrategia; puede fallar con personas fuera de la distribución \\
Value filter & Rechazar mensajes estratégicamente dañinos & Compara counterfactual expected value; ignora las consecuencias no modeladas \\
Closed loop & Hacer que el resultado de la comunicación vuelva a la siguiente planeación & Reestima la belief tras cada evento; aumentan la complejidad y la latencia \\
\bottomrule
\end{longtable}

\subsection{Conclusiones finales}

Primero, la escala de entrenamiento y el cómputo en inferencia son dos ejes de capacidad distintos. Poker y Go muestran que un base model fuerte puede seguir lejos del desempeño alcanzable si carece de search. Segundo, lo esencial de un agent de lenguaje no es la generación, sino el ciclo causal cerrado entre lenguaje, belief y action. Tercero, los datos humanos son a la vez un prior de comportamiento y el vehículo de las convention sociales, pero no sustituyen la optimización del valor. Cuarto, si un sistema no puede evaluar cierto tipo de riesgo de largo plazo, debe reconocer honestamente ese punto ciego y restringir el action space, en lugar de encubrir la ausencia de evaluator con una salida fluida.

\begin{importantbox}{La idea que más vale conservar}
Dedicar más cómputo a las decisiones correctas, confiar el modelado del mundo, la selección de estrategia y la expresión lingüística a módulos con contratos claros, y hacer que las consecuencias de la salida regresen a la siguiente ronda de planeación: este es el principio de diseño común que recorre esta sesión, de poker search a Cicero.
\end{importantbox}

\subsection{Lecturas adicionales}

Todos los materiales siguientes se publicaron antes de la fecha de la clase y forman directamente la cadena de evidencia de esta sesión. Se recomienda leer primero el artículo del sistema Cicero, después el artículo de piKL para entender el planner y, por último, revisar la tradición de search de poker y AlphaGo.

\begin{itemize}
  \item \href{https://www.science.org/doi/full/10.1126/science.ade9097}{Bakhtin et al., Human-level play in the game of Diplomacy by combining language models with strategic reasoning}: arquitectura del sistema Cicero, resultados en línea, diálogo y limitaciones.
  \item \href{https://proceedings.mlr.press/v162/jacob22a.html}{Jacob et al., Modeling Strong and Human-Like Gameplay with KL-Regularized Search}: objetivo formal de piKL y human-policy modeling.
  \item \href{https://proceedings.neurips.cc/paper/2021/hash/95f2b84de5660ddf45c8a34933a2e66f-Abstract.html}{Anthony et al., Learning to Play Diplomacy with No Human Supervision}: No-Press Diplomacy, juego contra sí mismo y el problema de los equilibrios múltiples.
  \item \href{https://www.science.org/doi/10.1126/science.aao1733}{Brown and Sandholm, Superhuman AI for heads-up no-limit poker}: Libratus y la resolución segura de subjuegos.
  \item \href{https://www.science.org/doi/10.1126/science.aay2400}{Brown and Sandholm, Superhuman AI for multiplayer poker}: Pluribus y el póquer multijugador.
  \item \href{https://www.science.org/doi/10.1126/science.aar6404}{Silver et al., Mastering the game of Go without human knowledge}: AlphaGo Zero y la search ablation.
  \item \href{http://www.incompleteideas.net/IncIdeas/BitterLesson.html}{Richard Sutton, The Bitter Lesson}: argumento histórico a favor de learning y search escalables.
\end{itemize}

\subsection{Autoevaluación posterior a la clase}

\begin{importantbox}{Ocho preguntas de autoevaluación}
\begin{enumerate}
  \item ¿En qué se diferencia la exploitability de la tasa de victorias contra un oponente fijo?
  \item ¿Por qué la slide 2 no demuestra que search siempre equivalga a un múltiplo fijo de base-model scaling?
  \item ¿Por qué la regla de support de Diplomacy hace que el lenguaje intervenga directamente en la transición de estados?
  \item ¿Por qué el intent conditioning puede mejorar a la vez la controllability y la perplexity?
  \item ¿Cómo falla cada uno de los enfoques puros, self-play e imitation?
  \item ¿Qué ocurre en piKL cuando $\lambda$ es demasiado grande o demasiado pequeña?
  \item ¿Por qué un value-based filter sigue sin poder evaluar de forma confiable el costo de mentir a largo plazo?
  \item Para trasladar la arquitectura de Cicero a un agent de mundo abierto, ¿qué transition, reward y límites de permisos faltan todavía?
\end{enumerate}
\end{importantbox}

\end{document}
